\chapter{Nilai Eigen}

Dalam bab tentang Dekomposisi Nilai Singular, kita membahas faktor yang
digunakan transformasi untuk mengubah panjang vektor dan menemukan faktor
perubahan panjang maksimum, yaitu nilai singular.
Dalam bab ini kita kembali menggambar bayangan lingkaran satuan untuk
membahas efek geometris lain, yaitu besar sudut pemutaran vektor.


\section{Pemutaran}
Ingat bahwa suatu pemetaan linear memberikan efek yang sama pada semua
vektor yang terletak pada garis yang sama melalui titik asal.
Jadi, kita dapat melihat aksinya dengan menggambar hasil pemetaan satu titik
pada setiap garis tersebut menggunakan rutin
\inlinecode{plot_circle_action}.\footnote{%
  Seperti dalam bab-bab sebelumnya, gambar-gambar ini dibuat dengan beberapa
  opsi yang tidak ditampilkan.
  Untuk daftar lengkapnya, lihat kode sumber panduan ini.}
\begin{sagecommandline}
sage: load("plot_action.sage")  
sage: q = plot_circle_action(1,0,0,1) 
sage: q.set_axes_range(-3, 3, -2, 3) 
sage: q.save("graphics/eigen000a.pdf") 
sage: p = plot_circle_action(2,0,2,2) 
sage: p.set_axes_range(-3, 3, -2, 3) 
sage: p.save("graphics/eigen000b.pdf")
\end{sagecommandline}
\begin{sagesilent}
load("plot_action.sage")  
q = plot_circle_action(1,0,0,1) 
q.set_axes_range(-3, 3, -3, 3) 
q.save("graphics/eigen000a.pdf", figsize=2.65)
p = plot_circle_action(2,0,2,2) 
p.set_axes_range(-3, 3, -3, 3) 
p.save("graphics/eigen000b.pdf", ticks_integer=True)
\end{sagesilent}
\begin{equation*} \label{gr:firstgraphic}
  \vcenteredhbox{\includegraphics{graphics/eigen000a.pdf}}
  \qquad\mapsunder{\big (\begin{smallmatrix} 2 &0 \\ 2 &2 \end{smallmatrix}\big )}\qquad
  \vcenteredhbox{\includegraphics{graphics/eigen000b.pdf}}
  % \tag{$*$}
\end{equation*}
Di sebelah kiri, jika ditelusuri berlawanan arah jarum jam, kurva dimulai
dengan warna merah di $(x,y)=(1,0)$.
Pada gambar di sebelah kanan, kurva dimulai dengan warna merah di $(1,2)$.
Berikut adalah vektor sebelum dan sesudah untuk titik merah tersebut.
\begin{sagecommandline}
sage: load("plot_action.sage")  
sage: p = plot_before_after_action(2,0,2,2, [(1,0)], ['red']) 
sage: p.set_axes_range(-0.5, 2, -0.5, 2) 
sage: p.save("graphics/eigen001.pdf", ticks=([1,2],[1,2]))
\end{sagecommandline}
\begin{sagesilent}
load("plot_action.sage")  
p = plot_before_after_action(2,0,2,2, [(1,0)], ['red']) 
p.set_axes_range(-0.5, 2, -0.5, 2) 
p.save("graphics/eigen001.pdf", ticks=([1,2],[1,2]), figsize=2)
\end{sagesilent}
% The input $\binom{1}{0}$ is mapped to the output $\binom{2}{2}$. 
Aksi matriks
\begin{equation*}
  \begin{mat}
   2  &0  \\
   2  &2
  \end{mat}
  \colvec{1 \\ 0}
  =\colvec{2 \\ 2}
  \qquad
  \vcenteredhbox{\includegraphics{graphics/eigen001.pdf}}
\end{equation*}
adalah mengubah panjang sekaligus memutar vektor.
\begin{sagecommandline}
sage: v = vector(RR, [1,0])
sage: M = matrix(RR, [[2, 0], [2, 2]])
sage: w = M*v
sage: v.norm(), w.norm()  
sage: angle_in_rads = arccos(w*v/(w.norm()*v.norm())) 
sage: angle_in_rads 
\end{sagecommandline}
Namun, untuk vektor-vektor lain, pemetaan ini akan mengubah panjang dan
memutarnya dengan besar yang berbeda.
Dengan kata lain, aksinya tidak seragam; ketika kita bergerak menyusuri
setengah lingkaran atas dari $(1,0)$ ke $(-1,0)$, efek pemetaan berubah-ubah.
Berikut adalah efek transformasi pada titik
$(\cos(\pi/6),\sin(\pi/6))$.
\begin{sagecommandline}
sage: v = vector(RR, [cos(pi/6),sin(pi/6)])
sage: w = M*v
sage: v.norm(), w.norm()  
sage: angle_in_rads = arccos(w*v/(w.norm()*v.norm())) 
sage: angle_in_rads 
sage: load("plot_action.sage")  
sage: p = plot_before_after_action(2,0,2,2, [(1,0), (cos(pi/6),sin(pi/6))], rainbow(12)) 
sage: p.set_axes_range(-0.5, 2, -0.5, 2.65) 
sage: p.save("graphics/eigen001a.pdf", ticks=([1,2],[1,2]))
\end{sagecommandline}
\begin{sagesilent}
v = vector(RR, [cos(pi/6),sin(pi/6)])
w = M*v
v.norm(), w.norm() 
angle_in_rads = arccos(w*v/(w.norm()*v.norm())) 
angle_in_rads 
load("plot_action.sage")  
p = plot_before_after_action(2,0,2,2, [(1,0), (cos(pi/6),sin(pi/6))], rainbow(12)) 
p.set_axes_range(-0.5, 2, -0.5, 2.65) 
p.save("graphics/eigen001a.pdf", ticks=([1,2],[1,2]), figsize=2)
\end{sagesilent}
Gambar ini kembali memperlihatkan keadaan sebelum dan sesudah untuk masukan
merah~$(1,0)$, lalu menumpangkan efek yang sama untuk masukan
oranye~$(\cos(\pi/6),\sin(\pi/6))$.
\begin{equation*}
  \vcenteredhbox{\includegraphics{graphics/eigen001a.pdf}}
\end{equation*}
Keluaran oranye tidak diputar sejauh keluaran merah dari masukannya
(panjangnya juga sedikit lebih besar).

\Sage{} dapat menghitung sudut pemutaran vektor-vektor tersebut.
Di akhir bab ini terdapat kode sumber rutin
\inlinecode{plot_color_angles}. Rutin itu menerima entri-entri sebuah
transformasi, menerapkan pemetaan tersebut pada vektor-vektor di setengah
lingkaran, menghitung sudut pemutaran vektor, lalu menggambar grafik.
Grafik itu mewarnai vektor-vektor keluaran sehingga kaitan visualnya terlihat.
\begin{sagecommandline}
sage: p = plot_color_angles(2,0,2,2)
sage: p.set_axes_range(0,pi,0,pi)
sage: p.save("graphics/eigen003.pdf")
\end{sagecommandline}
\begin{sagesilent}
load("plot_action.sage")  
p = plot_color_angles(2,0,2,2)
p.set_axes_range(0,pi,0,pi)
p.save("graphics/eigen003.pdf", figsize=3, tick_formatter=[pi,pi])
\end{sagesilent}
Transformasi untuk grafik ini adalah pemetaan yang, terhadap basis standar,
direpresentasikan oleh
$\left(\begin{smallmatrix}  2 &0 \\ 2 &2\end{smallmatrix}\right)$.
\begin{center}
  \vcenteredhbox{\includegraphics{graphics/eigen003.pdf}}
\end{center}
Titiknya yang paling menarik adalah ketika sudut keluaran bernilai $0$, yaitu
tempat grafik memotong sumbu horizontal.
Hal ini terjadi pada sudut masukan~$\pi/2$.
Vektor masukan ini sama sekali tidak diputar oleh transformasi\Dash panjangnya
diubah, tetapi arahnya tidak diputar.
Pada diagram masukan/keluaran di halaman~\pageref{gr:firstgraphic}, vektor itu
adalah vektor masukan hijau yang terletak pada sumbu~$y$; keluaran hijau yang
bersesuaian juga terletak pada sumbu~$y$.




\subsection{Matriks umum}
Kita dapat melakukan analisis yang sama untuk matriks umum $\nbyn{2}$ yang
biasa kita gunakan.
\begin{sagecommandline}
sage: q = plot_circle_action(1,0,0,1) 
sage: q.set_axes_range(-2, 2, -4, 5) 
sage: q.save("graphics/eigen100a.pdf")
sage: p = plot_circle_action(1,2,3,4) 
sage: p.set_axes_range(-2, 2, -4, 5) 
sage: p.save("graphics/eigen100b.pdf")
\end{sagecommandline}
\begin{sagesilent}
load("plot_action.sage")
q = plot_circle_action(1,0,0,1) 
q.set_axes_range(-2, 2, -4, 5) 
q.save("graphics/eigen100a.pdf", ticks_integer=True)
p = plot_circle_action(1,2,3,4) 
p.set_axes_range(-2, 2, -4, 5) 
p.save("graphics/eigen100b.pdf", ticks_integer=True)
\end{sagesilent}
\begin{equation*}
  \vcenteredhbox{\includegraphics{graphics/eigen100a.pdf}}
  \quad\mapsunder{\big (\begin{smallmatrix} 1 &2 \\ 3 &4 \end{smallmatrix}\big )}\quad
  \vcenteredhbox{\includegraphics{graphics/eigen100b.pdf}}
\end{equation*}
Grafik ini memperlihatkan sudut antara setiap vektor masukan dan vektor
keluaran yang bersesuaian.
% was:
% angles = find_angles(1,2,3,4,50)
% ticks = ([0,pi/4,pi/2,3*pi/4,pi], [0,pi/2,pi])
% p = scatter_plot(angles, marker='.', ticks=ticks)
% p.save("graphics/eigen101.pdf", figsize=2.5, tick_formatter=pi)
\begin{sagecommandline}
sage: p = plot_color_angles(1,2,3,4)
sage: p.set_axes_range(0,pi,-1*pi,pi)
sage: p.save("graphics/eigen101.pdf", tick_formatter=[pi,pi])
\end{sagecommandline}
\begin{sagesilent}
load("plot_action.sage")  
p = plot_color_angles(1,2,3,4)
p.set_axes_range(0,pi,-1*pi,pi)
p.save("graphics/eigen101.pdf", figsize=3, tick_formatter=[pi,pi])
\end{sagesilent}
\begin{equation*}
  \vcenteredhbox{\includegraphics{graphics/eigen101.pdf}}
  \tag{$*$}
\end{equation*}
Grafik ini mempunyai dua titik menarik: ketika $y=0$, dan ketika
$y=\pi$ (serta $y=-\pi$).
Di kedua tempat tersebut, panjang vektor berubah tetapi garis yang direntangnya
tidak berubah.
Pada kasus kedua, ketika vektor diputar dengan sudut~$\pi$ (atau $-\pi$),
vektor masukan dan keluaran terletak pada garis yang sama melalui titik asal,
tetapi transformasi mengubah skalanya dengan bilangan negatif.

Vektor yang oleh suatu transformasi~$t$ hanya dikenai perubahan skala disebut
\definend{vektor eigen} untuk~$t$ (vektor nol tidak kita izinkan sebagai
vektor eigen karena secara trivial dikenai perubahan skala oleh setiap
transformasi).
Faktor skala yang mengalikan vektor eigen disebut \definend{nilai eigen} yang
bersesuaian.
Untuk suatu nilai eigen~$\lambda$, himpunan vektor yang terkait dengan faktor
perubahan skala tersebut disebut \definend{ruang eigen}.
\Sage{} dapat menghitung ruang-ruang eigen sebuah matriks.\footnote{%
  Baik perkalian oleh vektor dilakukan dari kanan,~$M\vec{v}$, maupun dari
  kiri, $\vec{v}\,M$, sebuah matriks memiliki nilai-nilai eigen yang sama.
  Namun, vektor-vektor eigennya mungkin berbeda.
  Operasi \protect\Sage{} \protect\inlinecode{eigenvectors_right} menangani
  kasus standar buku ini, $M\vec{v}=\lambda\vec{v}$, sedangkan
  \protect\inlinecode{eigenvectors_left} menangani kasus yang lain.}

\begin{sagecommandline}
sage: M = matrix(RDF, [[1, 2], [3, 4]])
sage: evs = M.eigenvectors_right()
sage: evs
sage: evs[0] 
sage: v0 = vector(RDF, evs[0][1][0])
sage: v0.norm()
sage: evs[1]
sage: v1 = vector(RDF, evs[1][1][0])
sage: v1.norm()
\end{sagecommandline}
Hasil itu memperlihatkan daftar dengan dua elemen, masing-masing untuk satu
nilai eigen.
Elemen pertama, \inlinecode{evs[0]}, berkaitan dengan nilai eigen
$\lambda_1\approx -0.37$ dan menyatakan bahwa ruang eigen yang bersesuaian
mempunyai basis yang terdiri atas satu vektor satuan dengan titik ujung kira-kira
$(-0.82, 0.57)$.\footnote{%
  Angka~$1$ di bagian akhir adalah \protect\definend{multiplisitas geometris}
  nilai eigen ini, yaitu dimensi ruang nol
  $t-\lambda_1\cdot\identity$.
  Informasi ini tidak akan kita gunakan.}
Elemen kedua, \inlinecode{evs[1]}, berkaitan dengan
$\lambda_2\approx 5.37$ dan menyatakan bahwa ruang eigennya mempunyai basis
yang terdiri atas satu vektor dengan titik ujung kira-kira
$(-0.42, -0.91)$.

\Sage{} dapat menunjukkan letak masing-masing vektor tersebut pada grafik
berlabel~($*$).
\begin{sagecommandline}
sage: M = matrix(RDF, [[1, 2], [3, 4]])
sage: evs = M.eigenvectors_right()
sage: v = vector(RDF, evs[0][1][0])
sage: angle_v = atan2(v[1], v[0]) 
sage: (angle_v/pi).n(digits=4) 
\end{sagecommandline}
(Ingat bahwa fungsi \inlinecode{n()} memberikan nilai numerik argumennya.)
Jadi, sudutnya kira-kira $(4/5)\pi$. Dengan demikian, ruang eigen yang
dicantumkan pertama berkaitan dengan titik menarik di sebelah kanan pada
($*$), yaitu vektor yang diputar dengan sudut $\pi$.
Hal ini selaras dengan kenyataan bahwa nilai eigennya negatif, sebab keduanya
menyatakan bahwa aksi transformasi dari domain ke kodomain membalik arah
vektor.

\Sage{} dapat menggambar keadaan sebelum dan sesudah untuk kedua vektor eigen.
\begin{sagecommandline}
sage: M = matrix(RDF, [[1, 2], [3, 4]])
sage: evs = M.eigenvectors_right()  
sage: p1 = plot_before_after_action(1,2,3,4, [evs[0][1][0]], ['red']) 
sage: p2 = plot_before_after_action(1,2,3,4, [evs[1][1][0]], ['blue']) 
sage: p = p1+p2+circle((0,0), 1, edgecolor='black')
sage: p.set_axes_range(-3, 1, -5, 1) 
sage: p.save("graphics/eigen102.pdf")
\end{sagecommandline}
\begin{sagesilent}
load("plot_action.sage")  
M = matrix(RDF, [[1, 2], [3, 4]])
evs = M.eigenvectors_right()  
p1 = plot_before_after_action(1,2,3,4, [evs[0][1][0]], ['red']) 
p2 = plot_before_after_action(1,2,3,4, [evs[1][1][0]], ['blue'])
c = circle((0,0), 1, edgecolor='black') 
p = p1+p2+c
p.set_axes_range(-3, 1, -5, 1) 
p.save("graphics/eigen102.pdf", ticks_integer=True, aspect_ratio=1, figsize=4.5)
\end{sagesilent}
\noindent
Gambar ini memuat dua pasang vektor sebelum dan sesudah, satu pasang berwarna
biru dan yang lain berwarna merah (setiap vektor sebelum adalah vektor satuan).
Sekali lagi kita melihat bahwa vektor sesudah merupakan hasil perkalian skalar
vektor sebelum: pada kasus biru dengan faktor positif
$\lambda_2\approx 5.37$, dan pada kasus merah dengan faktor negatif
$\lambda_1\approx -0.37$.
\begin{equation*}
  \vcenteredhbox{\includegraphics{graphics/eigen102.pdf}}
\end{equation*}



\subsection{Rotasi bidang}
Kita mengetahui matriks $\nbyn{2}$ yang memutar semua vektor berlawanan arah
jarum jam.
\begin{equation*}
  \begin{mat}
    \cos\theta  &-\sin\theta  \\
    \sin\theta  &\cos\theta
  \end{mat}
\end{equation*}
Berikut adalah aksinya pada setengah lingkaran atas ketika $\theta=\pi/3$.
\begin{sagecommandline}
sage: load("plot_action.sage")
sage: q = plot_circle_action(1,0,0,1) 
sage: q.set_axes_range(-2, 2, -2, 2) 
sage: q.save("graphics/eigen200a.pdf")
sage: p = plot_circle_action(cos(pi/3),-sin(pi/3),sin(pi/3),cos(pi/3)) 
sage: p.set_axes_range(-2, 2, -2, 2) 
sage: p.save("graphics/eigen200b.pdf")
\end{sagecommandline}
\begin{sagesilent}
load("plot_action.sage")
q = plot_circle_action(1,0,0,1) 
q.set_axes_range(-1.25, 1.25, -1.25, 1.25) 
q.save("graphics/eigen200a.pdf", figsize=2.5)
p = plot_circle_action(cos(pi/3),-sin(pi/3),sin(pi/3),cos(pi/3))
p.set_axes_range(-1.25, 1.25, -1.25, 1.25) 
p.save("graphics/eigen200b.pdf", figsize=2.5)
\end{sagesilent}
\begin{equation*}
  \vcenteredhbox{\includegraphics{graphics/eigen200a.pdf}}
  \quad\mapsunder{\big (\begin{smallmatrix} \cos(\pi/3) &-\sin(\pi/3) \\ \sin(\pi/3) &\cos(\pi/3) \end{smallmatrix}\big )}\quad
  \vcenteredhbox{\includegraphics{graphics/eigen200b.pdf}}
\end{equation*}
Pemetaan ini memutar setiap vektor masukan sebesar $\pi/3$~radian, sehingga
grafik sudut masukan/keluarannya tidak mempunyai titik yang menarik.
\begin{sagecommandline}
sage: load("plot_action.sage")  
sage: p = plot_color_angles(cos(pi/3),-sin(pi/3),sin(pi/3),cos(pi/3))
sage: p.set_axes_range(0,pi,-0.25*pi,pi)
sage: p.save("graphics/eigen201.pdf", tick_formatter=[pi,pi])
\end{sagecommandline}
\begin{sagesilent}
load("plot_action.sage")  
p = plot_color_angles(cos(pi/3),-sin(pi/3),sin(pi/3),cos(pi/3))
p.set_axes_range(0,pi,-0.15*pi,pi)
p.save("graphics/eigen201.pdf", figsize=2.5, tick_formatter=[pi,pi])
\end{sagesilent}
\begin{equation*}
  \vcenteredhbox{\includegraphics{graphics/eigen201.pdf}}
\end{equation*}
Transformasi ini tidak mempunyai nilai eigen real.


% \subsection{Points}
% Another natural picture to draw is the action of linear maps on 
% some points.

% We start with the action of the linear transformation $\map{t}{\Re^2}{\Re^2}$
% represented with respect to the standard bases by our generic matrix with
% entries $1$, $2$, $3$, and $4$.
% \begin{sagecommandline}
% sage: load("plot_action.sage")
% sage: plot.options['figsize']=4.5
% sage: g = point_grid(3, 3)
% sage: p = plot_point_action(1, 2, 3, 4, g)
% sage: p.save("sagecommandline/plot_action110.pdf")  
% \end{sagecommandline}
% \begin{center}
%   \includegraphics{"sagecommandline/plot_action110.pdf"}
% \end{center}
% This picture shows that the eigenspace associated with the eigenvector
% having the largest absolute value is attractive.
% This is the \textit{power method}:

% \begin{sagecommandline}
% sage: load("plot_action.sage")
% sage: plot.options['figsize']=4.5
% sage: g = point_grid(5, 5)
% sage: p = plot_point_action(2, 2, 0, 2, g)
% sage: p.save("sagecommandline/plot_action111.pdf")  
% \end{sagecommandline}
% \begin{center}
%   \includegraphics{"sagecommandline/plot_action111.pdf"}
% \end{center}




\section{Polinomial matriks}

\Sage{} dapat mencari polinomial karakteristik dan polinomial minimal sebuah
matriks.
\begin{sagecommandline}
sage: M =  matrix(RDF, [[1, 2], [3, 4]]) 
sage: poly = M.charpoly()
sage: poly.factor()
sage: poly.roots()
\end{sagecommandline}
Tentu saja, nilai-nilai eigennya adalah akar-akar polinomial tersebut.

Ingat bahwa polinomial karakteristik dan polinomial minimal hanya dapat
berbeda jika polinomial karakteristik mempunyai akar berulang.
\begin{sagecommandline}
sage: M =  matrix(RDF, [[2, 2, 3], [0, 4, -1], [0, 0, 2]]) 
sage: M.charpoly()
sage: M.minpoly()
sage: M.charpoly().factor()
sage: M.minpoly().factor()
\end{sagecommandline}
Di sini \Sage{} kesulitan menentukan apakah $2$ merupakan akar berulang karena
masalah numerik titik mengambang.
Perhitungan eksak menggunakan bilangan rasional menyelesaikan persoalan ini.
\begin{sagecommandline}
sage: M =  matrix(QQ, [[2, 2, 3], [0, 4, -1], [0, 0, 2]]) 
sage: M.charpoly()
sage: M.minpoly()
sage: M.charpoly().factor()
sage: M.minpoly().factor()
\end{sagecommandline}

Matriks rotasi tidak mempunyai nilai eigen real, tetapi mempunyai nilai eigen
kompleks.
\begin{sagecommandline}
sage: R = matrix(RDF, [[cos(pi/3), -sin(pi/3)], [sin(pi/3), cos(pi/3)]])
sage: R.minpoly()
sage: R.eigenvalues()
sage: R.eigenvectors_right()
\end{sagecommandline}




\section{Diagonalisasi dan bentuk Jordan}

\Sage{} dapat memberi tahu kita apakah dua matriks saling serupa.
\begin{sagecommandline}
sage: S =  matrix(QQ, [[2, -3], [1, -1]]) 
sage: T =  matrix(QQ, [[0, -1], [1,  1]]) 
sage: S.is_similar(T)
sage: U =  matrix(QQ, [[1, 2], [3,  4]]) 
sage: S.is_similar(U)
\end{sagecommandline}
\noindent
% We can even get a transformation.
% \begin{sagecommandline}
% sage: S =  matrix(QQ, [[2, -3], [1, -1]]) 
% sage: T =  matrix(QQ, [[0, -1], [1,  1]]) 
% sage: S.is_similar(T, transformation=True)
% \end{sagecommandline}

Kita dapat menentukan apakah suatu matriks dapat didiagonalkan.
\begin{sagecommandline}
sage: M =  matrix(QQ, [[4, -2], [1, 1]])  
sage: M.is_diagonalizable()
\end{sagecommandline}
\noindent
Untuk mendiagonalkan matriks, ubahlah matriks itu ke bentuk Jordan.
(Sekalipun Anda belum mempelajari bentuk matriks ini, perintahnya sederhana.)
\begin{sagecommandline}
sage: M =  matrix(QQ, [[2, -2, 2], [0, 1, 1], [-4, 8, 3]])  
sage: M.jordan_form()
\end{sagecommandline}
Perhatikan garis \inlinecode{-+-+-} yang membagi matriks menjadi blok-blok
komponen horizontalnya.
Demikian pula, terdapat garis-garis vertikal.

\Sage{} juga dapat mencari matriks transformasi, yaitu matriks nonsingular~$P$
yang mengonversi antara matriks~$M$ yang diberikan dan matriks~$D$ dalam bentuk
yang diinginkan, dengan $M=PDP^{-1}$.
\begin{sagecommandline}[d,0,1]
sage: M =  matrix(QQ, [[2, -2, 2], [0, 1, 1], [-4, 8, 3]])  
sage: JF, P = M.jordan_form(transformation=True)
sage: JF
sage: P
sage: P^(-1)
sage: P*JF*P^(-1)
\end{sagecommandline}

% Even a non-diagonalizable matrix can be put in Jordan form.
% \begin{sagecommandline}
% sage: M =  matrix(QQ, [[2, -1], [1, 4]])  
% sage: M.jordan_form()
% \end{sagecommandline}



\section{Kode sumber plot\_color\_angles}
Rutin ini mengumpulkan objek-objek grafis dan mengembalikan plot dari daftar
objek tersebut.
Secara bawaan, rutin ini menggambar seribu titik.
\lstinputlisting[firstline=252, lastline=262]{plot_action.sage}

\section{Kode sumber color\_angles\_list}
Rutin ini mencari sudut yang dihasilkan oleh efek transformasi pada
vektor-vektor, lalu menggambar plot sebar titik-titik tersebut.
Setiap titik digambar dengan warna vektor masukannya.
Konstanta \inlinecode{TICKS} menentukan letak label pada sumbu.
\lstinputlisting[firstline=230, lastline=250]{plot_action.sage}

\section{Kode sumber find\_angles}
Rutin ini menggunakan rumus sudut antara dua vektor yang selalu menghasilkan
nilai positif; dengan kata lain, sudutnya tidak berorientasi.
Hal itu sesuai dengan tujuan kita di sini, yaitu menggunakan grafik untuk
memperkirakan tempat-tempat ketika aksi matriks tidak memutar vektor.
\lstinputlisting[firstline=200, lastline=227]{plot_action.sage}


\section{Kode sumber plot\_before\_after\_action}
Satu-satunya hal yang mungkin tak terduga dalam rutin ini dan rutin
pembantunya adalah: jika bayangan vektor tidak terlalu jauh,
\lstinputlisting[firstline=105, lastline=105]{plot_action.sage}
maka rutin pembantu tidak menampilkan panah, melainkan lingkaran.
\lstinputlisting[firstline=155, lastline=195]{plot_action.sage}
\endinput

PERLU DITINJAU
